\sectionstart
\section{总体架构与四引擎协作}
\label{sec:architecture}
\subsection{分层结构与信任边界}
系统采用同源网页、Python 标准库 HTTP 服务与 SQLite 数据库。页面负责输入、展示、选择和独立确认；网关解析会话及请求；业务服务读取事实、计算约束并执行状态变更。模型返回值与外部网页内容均作为待检查的数据，不因来源于模型而获得授权。

\begin{figure}[htbp]\centering
\caption{系统分层与对外连接边界}
\begin{tikzpicture}[node distance=7mm and 6mm]
\node[box,text width=12cm] (ui) {同源用户界面\quad 首页、工作台、财务体检、生活规划、模型中心};
\node[box,text width=12cm,below=of ui] (http) {HTTP 接入\quad 路由与请求校验、用户会话、读取同意};
\node[box,text width=5.5cm,below left=9mm and -5.65cm of http] (model) {受限智能编排\\意图、澄清、候选方案、只读回答};
\node[box,text width=5.5cm,right=6mm of model] (core) {确定性业务内核\\现金计算、版本复核、授权与记账};
\node[box,text width=5.5cm,below=of model] (ext) {外部服务\\模型与搜索、SMTP、美团连接器};
\node[box,text width=5.5cm,below=of core] (db) {SQLite 与独立运行状态\\合成账本、回执、审计、用户分区};
\draw[arr] (ui)--(http);\draw[arr] (http.south)-|(model.north);\draw[arr] (http.south)-|(core.north);
\draw[arr] (model)--node[above,font=\scriptsize]{校验候选}(core);
\draw[arr] (model)--(ext);\draw[arr] (core)--(db);
\end{tikzpicture}
\tabnote{外部连接不共享银行账本执行权。美团真实支付在用户与美团之间完成，订单状态与合成银行交易分别记录。}
\end{figure}

后端以模块而非独立模型进程组织职责：\code{server.py} 提供接入，\code{service.py} 管理账户、身份和账本，\code{conversation.py} 与 \code{agent_ops.py} 编排任务，\code{model.py} 适配模型，\code{life_planner.py} 执行有限期规划。四引擎是贯穿这些模块的逻辑职责，不要求四个大模型分别驻留，也不暗示存在未实现的自主协商网络。

\subsection{职责到工程模块的映射}
\begin{figure}[htbp]\centering
\caption{四种职责围绕同一份用户状态协作}
\begin{tikzpicture}[every node/.style={text=wpblue}]
\foreach \x/\g/\en/\cn in {0/知/ETHOS/理解处境,4.1/谋/METIS/形成安排,8.2/时/KAIROS/复核变化}{
  \node[draw=wpblue,circle,fill=wpcream,minimum size=15mm,line width=.8pt] (r\en) at (\x,0) {\Large\paperkaiti \g};
  \node[align=center,font=\small] at (\x,-1.22) {\en\\\cn};
}
\draw[-{Latex[length=2mm]},wpblue] (rETHOS.east)--(rMETIS.west);
\draw[-{Latex[length=2mm]},wpblue] (rMETIS.east)--(rKAIROS.west);
\node[draw=wprust,circle,fill=wpcream,minimum size=15mm,line width=.8pt,text=wprust] (guard) at (4.1,-3.1) {\Large\paperkaiti 守};
\node[align=center,font=\small,text=wprust] at (4.1,-4.3) {ARGOS 守诺\\授权、金额、版本与回执的跨层核验};
\draw[wprust,dashed] (0,-1.65)--(0,-1.85)--(8.2,-1.85)--(8.2,-1.65);
\draw[wprust,dashed] (guard.north)--(4.1,-1.85);
\end{tikzpicture}
\fignote{四者表示逻辑职责，并非四个独立大模型。守诺贯穿理解、规划和执行，不排在最后才介入。}
\end{figure}

\paragraph{ETHOS 知己}
以用户目标、期限、风险声明、当前现金与更正信息形成版本化画像，并解释当前任务侧重。实现位于画像规则及 \code{service.py}，长期声明与账本状态分别保存。

\paragraph{METIS 善谋}
结合账单、现金约束和允许的行动，生成候选方案与有序任务。实现分布在 \code{model.py}、\code{life_planner.py} 与 \code{agent_ops.py}；候选必须经过确定性核验。

\paragraph{KAIROS 择时}
检查收入事件、日期条件、有效期与版本变化，产出等待、失效和重新生成的状态。实现位于 \code{service.py}、\code{life_service.py} 与对话恢复流程。

\paragraph{ARGOS 守诺}
核对身份、授权、对象、金额、版本和回执，必要时拒绝推进。接入、业务服务和任务执行共同承担这项职责，最终状态变更受服务端约束。


ARGOS 的作用贯穿生成与执行：草案阶段检查可行性，确认阶段复核当前事实，执行阶段检查幂等与账本一致性。它拥有停止推进的业务效果，而不是在流程末尾补一句安全提示。KAIROS 的现有触发来自明确请求、演示事件与确认时检查；没有后台常驻事件执行器，条件满足后仍须重新核验与确认。

\subsection{画像到反馈的闭环}
用户首先声明目标、现金底线、用钱期限、流动性与风险意愿，并明确是否同意读取所需信息。画像层保留这些长期声明，账本层单独记录余额、账单、工资状态和版本。策略层据此约束候选空间，计划层形成可以逐项查看的步骤，业务层只执行当前有效且已经确认的动作。回执、用户更正和新事件成为后续判断的输入。

例如工资延期只改变短期收入时点，不等于用户同意提高风险或降低现金底线。现金路径重算后，原先可行的转账草案可能因账户版本变化而失效。系统应解释为何需要重看、保留已完成操作的回执，并为未执行部分重新准备草案；不把原来的授权延伸为新安排的空白授权。
